\documentclass[11pt,a4paper,twoside,titlepage=false,index=totoc,headings=normal]{scrreprt}
\usepackage[T1]{fontenc}
\usepackage{lmodern}
\usepackage{microtype}
\usepackage{amsmath,amssymb,amsthm}
\usepackage{enumitem}
\usepackage{xcolor}
\usepackage{tikz}
\usepackage{makeidx}
\usepackage[hidelinks]{hyperref}

% ---- Page: a wide outer margin ------------------------------------------------------
% The text block is about 12 cm wide. The outer margin (4.4 cm) holds \mnote remarks and
% leaves room for readers to write their own notes next to the text.
\usepackage[a4paper,twoside,includemp,inner=2.4cm,outer=1.6cm,
            marginparwidth=4.4cm,marginparsep=0.5cm,top=2.6cm,bottom=2.6cm]{geometry}
\newcommand*{\mnote}[1]{\marginpar{\footnotesize\raggedright\color{black!70}#1}}

% ---- The switch for solutions -------------------------------------------------------
% One line decides whether the solutions are typeset: \showsolutionstrue for your own
% copy, \showsolutionsfalse for the handout the students get.
\newif\ifshowsolutions
\showsolutionstrue

% ---- Boxes for definitions, theorems, examples and remarks --------------------------
\usepackage[most]{tcolorbox}   % "most" brings theorems, skins and breakable boxes
\definecolor{defcol}{HTML}{1F5F99}
\definecolor{thmcol}{HTML}{A32A2A}
\definecolor{excol}{HTML}{2E7D4F}
\definecolor{remcol}{HTML}{6B6B6B}
\tcbset{
  % Common look: a coloured bar on the left, a tinted background, no frame. The title
  % runs into the text ("Definition 1.2 (Name). The text ...") so boxes stay compact.
  notebox/.style={enhanced jigsaw, breakable, sharp corners, boxrule=0pt, frame hidden,
    left=9pt, right=6pt, top=4pt, bottom=4pt, before skip=9pt, after skip=9pt,
    fonttitle=\bfseries, separator sign={.}, description delimiters parenthesis,
    attach title to upper={\ }, after title={\ }},
  defstyle/.style={notebox, colback=defcol!6, coltitle=defcol, borderline west={2.5pt}{0pt}{defcol}},
  thmstyle/.style={notebox, colback=thmcol!6, coltitle=thmcol, borderline west={2.5pt}{0pt}{thmcol}},
  exstyle/.style={notebox, colback=excol!6, coltitle=excol, borderline west={2.5pt}{0pt}{excol}},
  remstyle/.style={notebox, colback=remcol!6, coltitle=remcol, borderline west={2.5pt}{0pt}{remcol}},
}
% Theorems, lemmas and corollaries share one counter, numbered by chapter: Theorem 1.3
% is followed by Definition 1.4. Use \begin{theorem}{Title}{label} and refer to it with
% Theorem~\ref{thm:label}. Definitions, lemmas and examples use the same pattern.
\newtcbtheorem[number within=chapter]{theorem}{Theorem}{thmstyle}{thm}
\newtcbtheorem[use counter from=theorem]{lemma}{Lemma}{thmstyle}{lem}
\newtcbtheorem[use counter from=theorem]{definition}{Definition}{defstyle}{def}
\newtcbtheorem[use counter from=theorem]{example}{Example}{exstyle}{ex}
\newtcolorbox{remark}{remstyle,title={Remark.}}   % remarks are not numbered

\renewcommand*{\qedsymbol}{$\blacksquare$}

% ---- Exercises with solutions -----------------------------------------------------------
\newcounter{exercise}[chapter]
\renewcommand*{\theexercise}{\thechapter.\arabic{exercise}}
\newenvironment{exercise}{\par\medskip\refstepcounter{exercise}\noindent\textbf{Exercise~\theexercise.}\ }{\par}
% The solution is only read from the file when the switch above is on; otherwise it
% vanishes without a trace. Put it directly after its exercise.
\NewEnviron{solution}{%
  \ifshowsolutions
    \begin{tcolorbox}[notebox,colback=black!4,coltitle=black!70,borderline west={2.5pt}{0pt}{black!40},
      title={Solution.},before skip=4pt]
      \BODY
    \end{tcolorbox}
  \fi}

% ---- Index: \defn{term} sets the term in italics and adds it to the index -------------
% (latexmk, which Lemmata uses, runs makeindex by itself; no extra setup is needed.)
\makeindex
\newcommand*{\defn}[1]{\emph{#1}\index{#1}}

% ---- Notation -------------------------------------------------------------------------
\newcommand*{\R}{\mathbb{R}}
\newcommand*{\N}{\mathbb{N}}
\newcommand*{\eps}{\varepsilon}
\DeclareMathOperator{\sgn}{sgn}

\subject{Lecture notes}
\title{Real Analysis I}
\subtitle{Course MATH 201 \textperiodcentered\ Autumn 2026}
\author{Prof.\ Alex Morgan\\ Department of Mathematics, Example University}
\date{Version of \today}

\begin{document}
\maketitle

\section*{How to read these notes}
Everything important sits in a box, so you can find it when you revise. Write your own
comments in the wide margin.

\begin{tcolorbox}[defstyle,title={Definition.}]
A new term is defined here and appears in the index.
\end{tcolorbox}
\begin{tcolorbox}[thmstyle,title={Theorem.}]
A result that is proved in the text; lemmas and corollaries look the same.
\end{tcolorbox}
\begin{tcolorbox}[exstyle,title={Example.}]
A worked case that makes the definition or theorem concrete.
\end{tcolorbox}
\begin{remark}
A side comment: intuition, a warning or a pointer to other topics.
\end{remark}

Exercises close each chapter. Solutions appear after the exercises only in the instructor
version of these notes.

% \tableofcontents would start a new page; \@starttoc prints the same list right here,
% so the first page is not left half empty. Use \tableofcontents on a long script.
\section*{Contents}
\makeatletter\@starttoc{toc}\makeatother

% ======================================================================================
\chapter{Sequences and limits}

\section{Convergence}
A \defn{sequence} of real numbers is a function $\N\to\R$, written $(a_n)_{n\ge 1}$.
We want to say precisely what it means for $a_n$ to get arbitrarily close to a number $a$.
\mnote{Think of $\eps$ as a challenge from an opponent, and of $N$ as your reply.}

\begin{definition}{Convergence}{conv}
  A sequence $(a_n)$ \defn{converges} to $a\in\R$ if for every $\eps>0$ there is an
  $N\in\N$ such that $|a_n-a|<\eps$ for all $n\ge N$. We then write $a_n\to a$ or
  $\lim_{n\to\infty}a_n=a$ and call $a$ the \defn{limit}.
\end{definition}

\begin{example}{The sequence $1/n$}{harmonic}
  Let $a_n=1/n$. Given $\eps>0$, choose $N>1/\eps$. For all $n\ge N$ we have
  $|a_n-0|=1/n\le 1/N<\eps$, so $a_n\to0$.
\end{example}

\begin{figure}[ht]
  \centering
  \begin{tikzpicture}[x=0.7cm,y=2.6cm]
    \fill[blue!10] (0.4,-0.2) rectangle (12.6,0.2);
    \draw[dashed,blue!60] (0.4,0.2) -- (12.6,0.2) (0.4,-0.2) -- (12.6,-0.2);
    \draw[->] (0.4,0) -- (12.9,0) node[right] {$n$};
    \node[blue!70!black,left] at (0.4,0.2) {\small $+\varepsilon$};
    \node[blue!70!black,left] at (0.4,-0.2) {\small $-\varepsilon$};
    \foreach \n in {1,...,12}{\fill[\ifnum\n<6 black\else blue!70!black\fi] (\n,1/\n) circle (1.6pt);}
    \draw[<->,red!70!black] (6,-0.35) -- node[below] {\small $n\ge N$} (12,-0.35);
  \end{tikzpicture}
  \caption{The points $a_n=1/n$ stay inside the strip of width $2\eps$ around the limit $0$
  from $N=6$ on (here $\eps=0.2$).}
  \label{fig:tube}
\end{figure}

\begin{theorem}{Uniqueness of the limit}{unique}
  A convergent sequence has exactly one limit.
\end{theorem}
\begin{proof}
  Suppose $a_n\to a$ and $a_n\to b$ with $a\ne b$, and let $\eps=|a-b|/2>0$. For large
  $n$ both $|a_n-a|<\eps$ and $|a_n-b|<\eps$ hold, so by the triangle inequality
  $|a-b|\le|a-a_n|+|a_n-b|<2\eps=|a-b|$, a contradiction.
\end{proof}

\begin{remark}
  Proofs in analysis often follow the same pattern as above: choose $\eps$ cleverly in
  terms of what you want to contradict, then take $n$ large enough for all the finitely
  many conditions you need.
\end{remark}

\section{Monotone sequences}
A sequence is \defn{monotone increasing} if $a_n\le a_{n+1}$ for all $n$, and
\defn{bounded} if $|a_n|\le M$ for some $M$.

\begin{theorem}{Monotone convergence}{mct}
  Every bounded monotone increasing sequence converges, and its limit is
  $\sup\{a_n : n\in\N\}$.
\end{theorem}
\begin{proof}
  Let $s=\sup_n a_n$, which exists because the set is non-empty and bounded above. Given
  $\eps>0$, the number $s-\eps$ is not an upper bound, so $a_N>s-\eps$ for some $N$. For
  $n\ge N$ monotonicity gives $s-\eps<a_N\le a_n\le s$, hence $|a_n-s|<\eps$.
\end{proof}

\begin{example}{Recursive sequence}{recursive}
  Let $a_1=1$ and $a_{n+1}=\sqrt{2+a_n}$. By induction $a_n<2$ and $a_n<a_{n+1}$, so
  Theorem~\ref{thm:mct} gives a limit $a$ with
  $a=\sqrt{2+a}$, that is $a=2$.
\end{example}

\section*{Exercises}
\begin{exercise}
  Show directly from Definition~\ref{def:conv} that $a_n=\dfrac{2n+1}{n+3}$ converges to $2$.
\end{exercise}
\begin{solution}
  We have $\left|\frac{2n+1}{n+3}-2\right|=\frac{5}{n+3}<\frac5n$. Given $\eps>0$, choose
  $N>5/\eps$; then $n\ge N$ implies $|a_n-2|<\eps$.
\end{solution}

\begin{exercise}
  Give an example of a bounded sequence that does not converge. Which hypothesis of
  Theorem~\ref{thm:mct} fails?
\end{exercise}
\begin{solution}
  $a_n=(-1)^n$ is bounded and does not converge; it is not monotone.
\end{solution}

% ======================================================================================
\chapter{Continuous functions}

\section{Continuity}
\begin{definition}{Continuity}{cont}
  A function $f:D\to\R$ with $D\subseteq\R$ is \defn{continuous} at $x_0\in D$ if for every
  $\eps>0$ there is a $\delta>0$ such that $|f(x)-f(x_0)|<\eps$ for all $x\in D$ with
  $|x-x_0|<\delta$. It is continuous if it is continuous at every point of $D$.
\end{definition}

\begin{lemma}{Sequence criterion}{seq}
  $f$ is continuous at $x_0$ if and only if $f(x_n)\to f(x_0)$ for every sequence $x_n\to x_0$
  in $D$.
\end{lemma}
\begin{proof}
  Omitted here. In your own notes, write the proof in the box above or the proof
  environment below it; long proofs break across pages without trouble.
\end{proof}

\begin{example}{A continuous and a discontinuous function}{signum}
  The function $x\mapsto x^2$ is continuous on $\R$. The sign function $\sgn$ is not
  continuous at $0$: the sequence $x_n=1/n\to0$ gives $\sgn(x_n)=1\not\to\sgn(0)=0$.
\end{example}

\section{The intermediate value theorem}
\begin{theorem}{Intermediate value theorem}{ivt}
  Let $f:[a,b]\to\R$ be continuous with $f(a)<0<f(b)$. Then $f(c)=0$ for some $c\in(a,b)$.
\end{theorem}
\begin{proof}
  Let $c=\sup\{x\in[a,b] : f(x)<0\}$. If $f(c)<0$, continuity would give points to the
  right of $c$ with $f<0$; if $f(c)>0$, it would give points to the left of $c$ with
  $f>0$, so that $c$ would not be the least upper bound. Hence $f(c)=0$.
\end{proof}

\begin{remark}
  The theorem fails over $\mathbb{Q}$: $x^2-2$ changes sign on $[0,2]$ but has no rational
  root. The completeness of $\R$ is exactly what the proof uses.
\end{remark}

\section*{Exercises}
\begin{exercise}
  Show that every polynomial of odd degree has a real root.
\end{exercise}
\begin{solution}
  Divide by the leading coefficient. For large $|x|$ the polynomial has the sign of
  $x^n$, which differs at $+R$ and $-R$ when $n$ is odd. Apply Theorem~\ref{thm:ivt}.
\end{solution}

% ======================================================================================
\printindex
\end{document}
